% A paper that exercises every field the LaTeX plugin extracts: a class
% with options, several packages, a sectioning outline, environments,
% labels, references, citations and a bibliography.
\documentclass[11pt,a4paper,twoside]{article}

\usepackage[utf8]{inputenc}
\usepackage{amsmath,amssymb}
\usepackage{graphicx}
\usepackage{hyperref}
\usepackage{booktabs}

\title{Reading a Working Copy Without Running a Command}
\author{Example Engineering}
\date{September 2026}

\begin{document}
\maketitle

\begin{abstract}
A checkout's index records enough to answer whether its tracked files have
changed, without invoking the source control system at all.
\end{abstract}

\section{Introduction}
\label{sec:intro}

The question a reader has is whether there is uncommitted work. See
Section~\ref{sec:method} for how we answer it, and \cite{torvalds2005} for
the format itself.

\section{Method}
\label{sec:method}

\subsection{The index}
\label{sec:index}

Each entry records a size and a modification time, so the comparison is
three-way:

\begin{equation}
\label{eq:compare}
s(f) =
\begin{cases}
  \text{changed} & \text{if } \lvert f \rvert \neq \lvert f_{\text{index}} \rvert \\
  \text{same}    & \text{if } t(f) = t_{\text{index}} \\
  h(f) = h_{\text{index}} & \text{otherwise}
\end{cases}
\end{equation}

\subsection{Padding}

Equation~\ref{eq:compare} is only reachable if the entries are walked
correctly, which needs the padding rule in Table~\ref{tab:padding}.

\begin{table}
\centering
\begin{tabular}{ll}
\toprule
Field & Bytes \\
\midrule
Fixed header & 62 \\
Path & variable \\
Padding & 1 to 8 \\
\bottomrule
\end{tabular}
\caption{An index entry.}
\label{tab:padding}
\end{table}

\begin{figure}
\centering
\includegraphics[width=0.6\textwidth]{padding}
\caption{Where a rounded-up offset lands.}
\label{fig:padding}
\end{figure}

\section{Results}
\label{sec:results}

Reading is enough. See Figure~\ref{fig:padding}.

\section*{Acknowledgements}

To whoever wrote the format down.

\bibliographystyle{plain}
\bibliography{references}

\end{document}
